En función de la utilidad que tendría el modelo para apoyar la evaluación de alineaciones y las decisiones de sustitución en baloncesto profesional. En este trabajo, dicha factibilidad no se refiere a la instalación de una aplicación en un equipo de la WNBA, sino a la posibilidad de generar recomendaciones de quintetos que sean interpretables, reproducibles y pertinentes para un proceso de análisis deportivo. La decisión final corresponde al cuerpo técnico, debido a que intervienen factores que pueden no estar representados en los registros históricos, como el estado físico de las jugadoras, las faltas acumuladas y las indicaciones tácticas de un encuentro.

\subsubsection*{Usuarios y utilidad prevista}

Los principales beneficiarios potenciales son los entrenadores, entrenadores asistentes y analistas de rendimiento, quienes deben valorar qué combinación de cinco jugadoras conviene utilizar frente a diferentes alineaciones rivales. El modelo permitiría comparar quintetos candidatos mediante estimaciones de rendimiento ofensivo y defensivo, en lugar de limitar la comparación a promedios individuales o a una sola estadística de equipo. En un escenario de análisis previo al encuentro, las estimaciones podrían utilizarse para estudiar alternativas de rotación; posteriormente, también podrían servir para contrastar las decisiones con lo observado en el partido.

Esta necesidad ha sido examinada en estudios empíricos. Chen et al. analizaron 3\,387 posesiones correspondientes a 23 partidos internacionales de baloncesto femenino y observaron diferencias entre las configuraciones y los estilos de alineación utilizados por distintas selecciones. Sus resultados relacionan el estudio de quintetos con la elección táctica de alineaciones y la selección de jugadoras \cite{Chen17122023}. Por otra parte, Grassetti et al. estudiaron datos \textit{play-by-play} de 240 partidos de la Euroliga 2018--2019 para estimar efectos individuales y colectivos mediante una extensión regularizada del \textit{adjusted plus-minus}. Los autores plantean que actualizar estas estimaciones a lo largo de la temporada puede aportar información útil para la gestión de quintetos \cite{grassetti2021lineupmanagement}. Estos trabajos constituyen antecedentes de la utilidad del análisis.

Las áreas de análisis de plantilla y reclutamiento podrían consultar los perfiles funcionales obtenidos en la primera fase. (Muniz y Flamand) aplicaron técnicas de agrupamiento a datos de la NBA para obtener características de juego más específicas que las posiciones tradicionales y señalaron su aplicación potencial en decisiones de selección, intercambios y contratación de jugadores \cite{muniz2022weightednetwork}. No obstante, las decisiones de la adquisición de jugadoras no forman parte del problema de optimización que se evaluará en el presente proyecto.

\subsubsection*{Situaciones documentadas que justifican el problema}

Una situación concreta ocurrió durante el quinto partido de las Finales de la WNBA de 2024, disputado entre New York Liberty y Minnesota Lynx el 20 de octubre. Ante las dificultades ofensivas de algunas de sus principales anotadoras, el cuerpo técnico de New York incrementó la participación de Nyara Sabally y utilizó una combinación de mayor tamaño junto con Breanna Stewart y Jonquel Jones. La decisión de ajustar la rotación y su contexto táctico fueron documentados por Dalzell; el encuentro terminó con victoria de New York por 67--62 después de tiempo extra, según el registro oficial de la WNBA \cite{dalzell2024libertyadjustment,wnba2024finalsgame5}. Este caso muestra que la selección de quintetos no se resuelve únicamente escogiendo a las cinco jugadoras con mejores estadísticas individuales: el contexto del rival y la complementariedad de la combinación también forman parte de la decisión. Se trata de un ejemplo de la necesidad deportiva, no de evidencia de que un algoritmo hubiera predicho ese resultado o de que la sustitución fuera su única causa.

La relevancia práctica del análisis de alineaciones también se observa en las fuentes disponibles para su consulta. La plataforma oficial \textit{WNBA Stats} dispone de una herramienta para examinar quintetos con filtros por temporada, minutos, situación de juego y rival, lo que confirma que el rendimiento de distintas combinaciones constituye una unidad de análisis reconocida por la liga \cite{wnbaStatsLineupsTool}. Asimismo, el sitio Basketball-Reference publica estadísticas históricas de alineaciones de cinco jugadoras, incluidos minutos compartidos y métricas por posesión \cite{basketballReferenceLiberty2024Lineups}. La existencia de estas herramientas demuestra que hay medios establecidos para consultar comparaciones de quintetos, aunque no permite afirmar, sin datos de uso, cuántos entrenadores las consultan o con qué frecuencia las utilizan.

Adicionalmente, la WNBA anunció en marzo de 2024 la incorporación de tecnología de seguimiento óptico tridimensional en sus arenas, orientada, entre otros fines, al análisis de jugadoras y equipos y al desarrollo de herramientas para estrategia deportiva \cite{wnba2024tracking}. Esta inversión de la liga constituye evidencia de interés institucional en el análisis avanzado del juego. Sin embargo, el proyecto no presupone acceso al sistema comercial de seguimiento; su desarrollo se limita a los datos públicos disponibles y a las variables que puedan reconstruirse de manera verificable.

\subsubsection*{Condiciones para su utilización}

Para que una recomendación tenga utilidad operativa, deberá identificar el quinteto rival considerado, las cinco jugadoras propuestas y una estimación comparable de su rendimiento, junto con información sobre la cantidad de observaciones que respalda el resultado. Asimismo, las alternativas deberán respetar la disponibilidad de las jugadoras y las restricciones de sustitución definidas en el modelo. Cuando existan pocos antecedentes de una combinación, la estimación deberá acompañarse de una indicación de incertidumbre; posteriormente se podrá hacer uso de algún método de imputación. 

La evaluación operativa se concentrará inicialmente en la preparación de escenarios previos o posteriores a los partidos, sin asumir capacidad de respuesta en tiempo real. Se comprobará si el procedimiento produce resultados interpretables y reproducibles, cuánto tarda en comparar los quintetos factibles y si la información resultante añade elementos de análisis frente a las métricas disponibles en las fuentes de consulta. Si se cuenta con personas con experiencia en entrenamiento o análisis de baloncesto, podrá evaluarse la claridad y pertinencia de las recomendaciones mediante escenarios estructurados. Tal valoración complementaría la validación cuantitativa, pero no sustituiría la medición del error predictivo y la calidad del ordenamiento.

En consecuencia, se considera que el proyecto tiene una \textbf{utilidad operativa potencial} para entrenadores y analistas deportivos, sustentada tanto en antecedentes de evaluación de quintetos como en casos reales de ajustes de alineación y en la existencia de herramientas oficiales de consulta estadística. La factibilidad operativa definitiva del modelo propuesto quedará condicionada a la calidad de los datos reconstruidos, el desempeño obtenido en la validación temporal, la interpretación de sus resultados y la posibilidad de ejecutar el procedimiento con los recursos disponibles.